\documentclass[11pt]{article}

% This file will be kept up-to-date at the following GitHub repository:
%
% https://github.com/automl-conf/LatexTemplate
%
% Please file any issues/bug reports, etc. you may have at:
%
% https://github.com/automl-conf/LatexTemplate/issues

\usepackage{microtype} % microtypography
\usepackage{booktabs}  % tables
\usepackage{url}  % urls
\usepackage{csquotes}

% AMS math
\usepackage{amsmath}
\usepackage{amsthm}


\usepackage{graphicx} % added

\usepackage{float}

% With no package options, the submission will be anonymized, the supplemental
% material will be visible, and line numbers will be added to the manuscript.
%
% To compile a non-anonymized camera-ready version, add the [final] option:
%
% \usepackage[final]{automl}
%
% To compile for the non-archival content track, add the [shortpaper] option:
%
% \usepackage[shortpaper]{automl}
% \usepackage[final,shortpaper]{automl}
%
% To compile a single-blind submission, e.g. for the ABCD track, add the
% [revealauthors] option:
%
\usepackage[revealauthors]{automl}
%
% To compile a non-anonymized pre-print version that can be uploaded to arXiv,
% add the [preprint] option:
%
% \usepackage[preprint]{automl}
%
% The [hidesupplement] option may be used to hide the supplementary material from any of the
% above commands, e.g.:
%
% \usepackage[hidesupplement,final,shortpaper]{automl}

% \usepackage[final]{automl}

% You may use any reference style as long as you are consistent throughout the
% document. As a default, we suggest author--year citations; for bibtex and
% natbib you may use:

\usepackage{natbib}
\bibliographystyle{apalike2}

% and for biber and biblatex you may use:

% \usepackage[%
%   backend=biber,
%   style=authoryear-comp,
%   sortcites=true,
%   natbib=true,
%   giveninits=true,
%   maxcitenames=2,
%   doi=false,
%   url=true,
%   isbn=false,
%   dashed=false
% ]{biblatex}
% \addbibresource{example.bib}


\title{Surrogate Neural Architecture Codesign Package (SNAC-Pack)}

% The syntax for adding an author is
%
% \author[i]{\nameemail{author name}{author email}}
%
% where i is an affiliation counter. Authors may have
% multiple affiliations; e.g.:
%
% \author[1,2]{\nameemail{Anonymous}{anonymous@example.com}}
%
% the list might continue:
% \author[2,3]{\nameemail{Author 2}{email2@example.com}}
% \author[3]{\nameemail{Author 3}{email3@example.com}}
%
% if you need to force a linebreak in the author list, prepend an \author entry
% with \\:
%
% \author[1,2,3]{\\\nameemail{Author 4}{email4@example.com}}
%
% Specify corresponding affiliations after authors, referring to counter used in
% \author:
%
% \affil[1]{Institution 1}
%
% the list might continue:
% \affil[2]{Institution 2}
% \affil[3]{Institution 3}
%
% this system may also be used to add additional footnotes, e.g. to indicate
% equal contributions. Here is an example:
%
% \author[1,$\ast$]{\nameemail{Author 1}{email1@example.com}}
% \author[1,$\ast$]{\nameemail{Author 2}{email2@example.com}}
% \author[2]{\nameemail{Author 3}{email3@example.com}}
%

\author[1]{\nameemail{Jason Weitz}{jdweitz@ucsd.edu}}
\author[2,1]{\nameemail{Dmitri Demler}{ddemler@ethz.ch}}
\author[3]{\nameemail{Benjamin Hawks}{bhawks@fnal.gov}}
\author[4]{\nameemail{Aaron Wang}{aaronw5@uic.edu}}
\author[3]{\nameemail{Nhan Tran}{ntran@fnal.gov}}
\author[1]{\nameemail{Javier Duarte}{jduarte@ucsd.edu}}

% \author{%
%   Jason Weitz\\
%   University of California San Diego\\
%   La Jolla, CA 92093, USA \\
%   \texttt{jdweitz@ucsd.edu} \\
%   \And
%   Dmitri Demler \\
%   University of California San Diego\\
%   La Jolla, CA 92093, USA\\
%   \texttt{ddemler@ucsd.edu} \\
%   \AND
%   Benjamin Hawks \\
%   Fermi National Accelerator Laboratory \\
%   Batavia, IL 60510, USA \\
%   \texttt{bhawks@fnal.gov} \\
%   \And
%   Nhan Tran \\
%   Fermi National Accelerator Laboratory \\
%   Batavia, IL 60510, USA \\
%   \texttt{ntran@fnal.gov} \\
%   \And
%   Javier Duarte \\
%   University of California San Diego\\
%   La Jolla, CA 92093, USA \\
%   \texttt{jduarte@ucsd.edu} \\
% }

\affil[1]{University of California San Diego}
\affil[2]{ETH Zurich}
\affil[3]{Fermi National Accelerator Laboratory }
\affil[4]{University of Illinois Chicago}

% define PDF metadata, aids in accessibility of the resulting PDF
\hypersetup{%
  pdfauthor={Jason Weitz, Dmitri Demler, Benjamin Hawks, Aaron Wang, Nhan Tran, Javier Duarte}, % will be reset to "Anonymous" unless the "final" package option is given
  pdftitle={Surrogate Neural Architecture Codesign Package (SNAC-Pack)},
  pdfsubject={AutoML Conference Submission},
  pdfkeywords={neural architecture search, FPGA, hardware-aware}
}

\begin{document}

\maketitle

\begin{abstract}
Neural architecture search (NAS) is a powerful approach for automating model design, but existing methods often optimize for accuracy alone or rely on proxy metrics such as bit operations (BOPs) that correlate poorly with hardware cost.
This gap is particularly large for FPGA deployment, where cost is dominated by a multi-dimensional budget of lookup tables, DSPs, flip-flops, BRAM, and latency.
We present the Surrogate Neural Architecture Codesign Package (SNAC-Pack), an open-source AutoML framework for hardware-aware neural architecture codesign and end-to-end FPGA deployment.
SNAC-Pack runs a multi-objective global search with Optuna and NSGA-II, loading trials to a shared SQLite store that enables parallel workers across compute nodes.
A hardware surrogate model outputs per-trial resource and latency estimates, avoiding the synthesis cost that would otherwise dominate the search loop.
A local search stage then applies quantization-aware training (QAT) together with iterative magnitude pruning in a combined compression loop, after which the final model is synthesized to FPGA firmware via the hls4ml Python library.
A YAML configuration and an optional agentic frontend let users run the pipeline on new datasets without modifying the framework.
We demonstrate SNAC-Pack on jet classification at the Large Hadron Collider and superconducting qubit readout, discovering compact architectures that match or exceed strong baselines on the task metric while reducing FPGA resource utilization.
\end{abstract}

\section{Introduction}
\label{sec:intro}

Machine learning models now achieve state-of-the-art performance across many domains, but deploying them in resource-constrained environments such as edge devices, embedded accelerators, and low-latency triggers remains a substantial engineering challenge.
FPGAs are a natural target in such settings because they offer microsecond-scale inference and reconfigurable parallelism, but they impose a multi-dimensional resource budget spanning lookup tables (LUTs), digital signal processors (DSPs), block RAM (BRAM), flip-flops (FFs), and clock cycles, all of which depend on model architecture, quantization, and pruning.
Neural architecture search (NAS) automates model design and can in principle navigate these trade-offs, but most existing methods either optimize for accuracy alone or rely on software-level proxies such as bit operations (BOPs) and floating-point operations (FLOPs) that correlate only loosely with post-synthesis hardware cost.
Closing the gap between the metrics optimized during search and the metrics that dictate deployment requires fast, accurate hardware feedback inside the search loop and a workflow that connects search to physical synthesis end-to-end.

% We introduce the Surrogate Neural Architecture Codesign Package (SNAC-Pack), an open-source automated machine learning (AutoML) framework for hardware-aware neural architecture codesign with end-to-end FPGA deployment.
% SNAC-Pack runs a multi-objective global search with Optuna~\citep{optuna_2019} and an NSGA-II~\citep{NSGA} search algorithm, where each trial samples an architecture, predicts its FPGA resource utilization and latency with a learned surrogate~\citep{Rahimifar_2025}, and trains the model to obtain accuracy.
% To enable practical search at scale, all trials are written to a shared SQLite store, allowing workers on independent compute nodes to execute trials in parallel against a common Pareto archive.
% A selected optimal candidate then enters a local search stage that applies iterative magnitude pruning and quantization-aware training~\citep{quantization_survey,modelcompression}, after which the final model is synthesized to FPGA firmware via the hls4ml Python library~\citep{Schulte:2025mai,hls4ml,fastml_hls4ml}.
% Users drive the full pipeline through a YAML configuration that specifies the dataset, search space, optimization objectives, and target hardware; an optional agentic frontend based on the Model Context Protocol allows for the same configuration options using natural-language with an MCP agent for non-expert users.

We introduce the Surrogate Neural Architecture Codesign Package (SNAC-Pack), an open-source AutoML framework for hardware-aware neural architecture codesign through FPGA deployment (Fig.~\ref{fig:pipeline_fig}; \S\ref{sec:method}).
Relative to Neural Architecture Codesign (NAC)~\citep{Weitz_2025}, SNAC-Pack scores global-search trials with learned surrogate estimates of utilization and latency~\citep{Rahimifar_2025} instead of relying on BOPs alone, while preserving multi-objective Optuna~\citep{optuna_2019} search with NSGA-II~\citep{NSGA}, local compression (quantization-aware training and pruning~\citep{quantization_survey,modelcompression}), and hls4ml synthesis~\citep{Schulte:2025mai,hls4ml,fastml_hls4ml}.
Users configure datasets, search spaces, objectives, and hardware in YAML; an optional Model Context Protocol front end exposes the same workflow to agentic clients.


We demonstrate SNAC-Pack on two scientific machine learning tasks with different deployment profiles: jet classification at the Large Hadron Collider~\citep{pierini_2020_3602260}, where sub-microsecond inference is required at the trigger level, and superconducting qubit readout~\citep{diguglielmo2025endtoendworkflowmachinelearningbased}, where readout fidelity must be balanced against strict hardware budgets.
In both cases, SNAC-Pack finds compact architectures that match or exceed baselines on the task metric while reducing post-synthesis FPGA resource utilization.
SNAC-Pack is open-source and publicly available~\citep{dmitri_demler_2026_19202843}.

% , and we compare directly against an otherwise-identical pipeline using BOPs as the efficiency objective to isolate the contribution of hardware-aware surrogate feedback. 

\begin{figure}
    \centering
    \includegraphics[width=0.95\linewidth]{figures/snac_figure_try5.pdf}
    % \caption{End-to-end SNAC-Pack pipeline.
    % The user configures the run via a YAML file or an optional MCP agent interface.
    % The Surrogate Codesign Hub runs multi-objective global search, where each trial samples an architecture, predicts FPGA resources and latency with a learned surrogate, evaluates accuracy, and updates the Pareto front.
    % A selected candidate undergoes local search (iterative magnitude pruning and quantization-aware training) before synthesis to FPGA firmware with hls4ml.}
    \caption{YAML (or optional MCP) drives global multi-objective search with surrogate hardware scores, local QAT and iterative magnitude pruning, and hls4ml synthesis.}
    \label{fig:pipeline_fig}
\end{figure}

\section{Related Work}
\label{sec:related}

% SNAC-Pack draws on three threads of prior work: neural architecture search, surrogate models for hardware estimation, and AutoML
% frameworks.
% It builds on Neural Architecture Codesign (NAC)~\citep{Weitz_2025}, a multi-stage NAS in which a global search produces a Pareto front of candidate architectures and a local search refines selected models through quantization-aware training~\citep{quantization_survey} and pruning~\citep{modelcompression} before synthesis with hls4ml.
% Its efficiency objective is BOPs, a software proxy that only loosely tracks FPGA cost.
% SNAC-Pack closes this gap by combining rule4ml~\citep{Rahimifar_2025}, a learned surrogate for FPGA resource utilization and latency, directly into the search loop.
% The remainder of this section splits these components within the broader literature: NAS and its multi-objective extensions, hardware-aware NAS, surrogate models for hardware estimation, and AutoML frameworks targeting deployment.

SNAC-Pack builds on Neural Architecture Codesign (NAC)~\citep{Weitz_2025}: a global search produces a Pareto front, local search refines candidates with quantization-aware training~\citep{quantization_survey} and pruning~\citep{modelcompression}, and hls4ml~\citep{hls4ml,fastml_hls4ml} performs synthesis.
NAC's efficiency objective is BOPs, which only loosely tracks FPGA cost; SNAC-Pack adds rule4ml~\citep{Rahimifar_2025} in the global loop to estimate utilization and latency before invoking full synthesis for every trial.
The following subsections review neural architecture search, hardware-aware NAS, surrogate estimation, and AutoML tooling in that order.

\subsection{Neural Architecture Search}
% Neural Architecture Search (NAS) automates the design of neural networks by exploring a space of candidate architectures under a chosen search strategy.
% \citet{zoph2016neural} established the modern formulation, training a reinforcement learning controller to produce architectures scored by validation accuracy but at a cost of thousands of GPU-days per search.
% Subsequent work reduced this cost dramatically through weight sharing~\citep{pham2018efficient} and continuous relaxation of the search space~\citep{liu2018darts}, but both introduce constraints that complicate hardware-aware deployment: shared-weight rankings can be biased, and gradient-based formulations do not naturally accommodate non-differentiable objectives such as FPGA resource utilization.

% Another work reframed NAS as an inherently multi-objective problem.
% The survey by \citet{9508774} catalogs the broad use of evolutionary computation for NAS, where populations of architectures are evolved via mutation and crossover and ranked across multiple objectives simultaneously.
% \citet{lu2019nsga} applied NSGA-II~\citep{NSGA} to jointly optimize classification error and computational complexity, producing Pareto fronts of architectures rather than single solutions.
% \citet{elsken2018efficient} similarly optimized for accuracy and resource cost, using network morphism to inherit weights across generations and offset training cost.
% \citet{sinha2024multi} extend NSGA-II with zero-cost proxies and explicit diversity preservation to broaden the resulting Pareto fronts.

% A common limitation across these multi-objective methods is the choice of efficiency metric.
% FLOPs, parameter counts, and BOPs are attractive because they are cheap to compute, but they do not reliably reflect the metrics that govern real deployment on reconfigurable hardware: LUTs, DSPs, BRAM, FFs, and clock cycles.
% Closing this gap motivates hardware-awareness.

Neural Architecture Search (NAS) automates architecture design under a search strategy.
\citet{zoph2016neural} introduced reinforcement-learning NAS that was accurate but extremely costly; weight sharing~\citep{pham2018efficient} and DARTS-style relaxation~\citep{liu2018darts} reduced cost but complicate hardware-aware use because supernet rankings can be biased and gradient-based encodings handle non-differentiable FPGA objectives poorly.
Multi-objective evolutionary NAS is surveyed by \citet{9508774}; \citet{lu2019nsga} applied NSGA-II~\citep{NSGA} to trade classification error against complexity and return Pareto fronts, \citet{elsken2018efficient} optimized accuracy and resource cost with network morphism across generations, and \citet{sinha2024multi} extended NSGA-II with zero-cost proxies and diversity preservation to widen Pareto coverage.
Despite this progress, efficiency is still commonly reported with FLOPs, parameter counts, or BOPs, which are cheap but misaligned with reconfigurable-hardware budgets (LUTs, DSPs, BRAM, FFs, clock cycles), motivating hardware-aware NAS.


\subsection{Hardware-Aware NAS}
% Hardware-aware NAS emerged in response to the observation that proxy metrics correlate poorly with real deployment cost.
% \citet{tan2019mnasnet} measured inference latency directly on mobile devices and incorporated it into the search reward, demonstrating that FLOPs alone routinely mispredict runtime.
% To avoid the cost of per-trial on-device measurement, \citet{wu2019fbnet} precomputed a latency lookup table over operators and treated their contributions as additive within a differentiable search.
% \citet{cai2018proxylessnas} learned a latency prediction model, enabling search to run directly on the target task and hardware without intermediate proxy datasets.

Hardware-aware NAS targets true deployment cost because FLOP-style proxies correlate poorly with realized performance.
\citet{tan2019mnasnet} incorporated measured mobile latency into the search reward; \citet{wu2019fbnet} avoided repeated on-device calls with additive operator latency tables inside differentiable search; \citet{cai2018proxylessnas} learned latency predictors so candidates could be scored on target hardware without intermediate proxy datasets.
Later work widened the design space: \citet{cai2019once} decoupled supernet training from evolutionary subnet selection using fast accuracy and latency predictors, \citet{wang2020apq} co-searched architecture with pruning and mixed precision, and \citet{benmeziane2023multi} trained rank-preserving surrogates aligned with Pareto fronts rather than absolute error.

% Subsequent work expanded the scope of hardware-aware search.
% \citet{cai2019once} decoupled training from architecture selection: a single supernet is trained once, after which evolutionary search over fast accuracy and latency predictors yields sub-networks for arbitrary hardware targets. \citet{wang2020apq} jointly searched over architecture, pruning ratios, and mixed-precision quantization policies, showing that optimizing these axes sequentially leaves substantial efficiency unrealized.
% \citet{benmeziane2023multi} further refined this approach by training surrogates with rank-preserving losses tuned to the Pareto front rather than to absolute prediction error.

% These methods overwhelmingly target mobile CPUs, GPUs, and edge accelerators, where latency is the dominant cost and can be modeled as an approximately additive scalar.
% FPGAs differ in kind: deployment is bound by a multi-dimensional budget of LUTs, DSPs, BRAM, and FFs, with nontrivial interactions between quantization, pruning, and resource allocation.
% \citet{jiang2019accuracy} proposed one of the earliest FPGA-aware NAS frameworks, bounding latency through an analytical abstraction over loop-unrolling and pipelining; \citet{li2021hw} later released a benchmark of FPGA and edge-device costs over standard NAS search spaces.
% Both lower the barrier to FPGA-aware research but rely on static or analytical models; closing the loop on dynamic search over arbitrary architectures motivates surrogate-based approaches.

Much of this line of work still centers mobile CPUs, GPUs, and edge accelerators, where scalar latency is a workable surrogate.
FPGAs instead bind LUT, DSP, BRAM, and FF utilization together with nontrivial quantization and pruning interactions; \citet{jiang2019accuracy} proposed an early FPGA-aware NAS framework with analytical latency bounds over loop unrolling and pipelining, and \citet{li2021hw} benchmarked FPGA and edge costs on standard NAS spaces, yet both lean on static or analytical models rather than closing the loop for arbitrary architectures, motivating surrogate-based search.

\subsection{Surrogate Models for Hardware Estimation}
% The hardware-aware methods above all rely on a model mapping each candidate architecture to its deployment cost.
% When that cost comes from full synthesis or on-device profiling, every evaluation takes minutes to hours, making it intractable inside a search loop over thousands of candidates.
% Surrogate models replace this expensive ground truth with a fast learned predictor.
% For edge CPUs, GPUs, and mobile accelerators, \citet{zhang2021nn} showed that fusion-aware kernel-level prediction is necessary to model runtime behavior under graph optimization, and \citet{akhauri2024latency} subsequently studied the transferability of such predictors across devices.
% \citet{dudziak2020brp} brought the surrogate paradigm into NAS itself through a GNN trained on binary relations between architectures, improving sample efficiency when labeled data is scarce.

Hardware-aware search still needs a per-candidate map to deployment cost, but synthesis or on-device profiling often costs minutes to hours per evaluation, so large search budgets rely on fast learned surrogates instead of repeated ground truth.
For edge CPUs, GPUs, and mobile accelerators, \citet{zhang2021nn} argued fusion-aware kernel-level prediction is needed under graph optimization, \citet{akhauri2024latency} studied transfer of such predictors across devices, and \citet{dudziak2020brp} embedded a GNN over binary architecture relations in NAS to improve sample efficiency when labels are scarce.

% These approaches treat deployment cost as an additive scalar, which breaks down on FPGAs where an architecture must simultaneously fit within budgets of LUTs, DSPs, FFs, and BRAM, and these resources interact non-trivially with quantization and pruning.
% Earlier FPGA-targeted predictors operated at the level of HLS code: \citet{10.1145/3240765.3264635} trained regression models on extracted code features to predict post-place-and-route performance without invoking HLS in the loop, and \citet{10.1145/3489517.3530408} extended this to a graph neural network predictor over a 40k-program benchmark of C/C++ kernels, outperforming HLS-based estimates on both resource and timing metrics.
% The rule4ml library~\citep{Rahimifar_2025} specializes this idea to neural networks by training per-resource regressors over a synthesized dataset of NN architectures, predicting LUT, DSP, FF, and BRAM utilization together with II and latency clock cycles before any HLS run.
% The wa-hls4ml benchmark and accompanying GNN and transformer-based surrogates~\citep{hawks2025wahls4mlbenchmarksurrogatemodels} extend this with predictors trained on over 680{,}000 hls4ml-synthesized networks spanning fully connected and convolutional architectures, providing a standardized evaluation suite. SNAC-Pack uses rule4ml with the wa-hls4ml GNN as a surrogate inside the search loop and reserves full hls4ml synthesis for final validation, treating the surrogate as a swappable component so future predictors can replace it without modifying the search interface.

Scalar latency surrogates align poorly with FPGAs, where LUT, DSP, FF, and BRAM utilization co-vary with quantization and pruning.
\citet{10.1145/3240765.3264635} trained regressors on HLS code features to predict post-place-and-route performance without invoking HLS in the loop; \citet{10.1145/3489517.3530408} extended this with a GNN over a 40k-program C/C++ benchmark, outperforming HLS-based estimates on resource and timing.
rule4ml~\citep{Rahimifar_2025} specializes the idea to NNs with per-resource regressors on synthesized architectures, forecasting utilization, II, and latency cycles before HLS; wa-hls4ml~\citep{hawks2025wahls4mlbenchmarksurrogatemodels} adds GNN and transformer baselines trained on over 680{,}000 hls4ml-synthesized fully connected and convolutional models as a standardized suite.
SNAC-Pack runs rule4ml with the wa-hls4ml GNN inside the search loop, reserves full hls4ml synthesis for final validation, and keeps the surrogate swappable so future predictors can plug in without changing the search interface.


\subsection{AutoML Frameworks and Tooling}
% The AutoML landscape has evolved from algorithm selection over fixed feature pipelines to fully automated deep learning.
% Auto-sklearn~\citep{feurer2015efficient} established the early paradigm by combining Bayesian optimization with meta-learning warm-starts and post-hoc ensembling over the scikit-learn library, targeting classical CPU-bound classifiers.
% As deep learning matured, frameworks shifted toward joint architecture and hyperparameter search: AutoKeras~\citep{jin2023autokeras} encapsulated NAS and hyperparameter tuning behind a high-level Keras API to broaden accessibility, while
% Auto-PyTorch~\citep{zimmer2021auto} jointly optimized architectures
% and training schedules through multi-fidelity meta-learning for
% tabular AutoDL.
% To address the cost of search itself, FLAML~\citep{wang2021flaml} introduced budget-adaptive randomized search and search-space decomposition, prioritizing efficient model discovery rather than asymptotic accuracy.
% Optuna~\citep{optuna_2019} provides a flexible programmatic search interface with distributed execution that has become a common backbone for modern AutoML pipelines.

% These frameworks largely treat deployment as a downstream concern: their objectives are software-level metrics such as validation accuracy or GPU throughput, and their backends end at a trained-model artifact.
% AutoGluon-Multimodal~\citep{tang2024autogluon} extends this line to multimodal foundation models with metrics such as inference frames-per-second on high-end GPUs, but these signals do not translate to the physical-resource budgets that govern reconfigurable hardware.
% hls4ml~\citep{hls4ml,fastml_hls4ml,10.1145/3801979} closes part of this gap on the deployment side by translating quantized Keras and PyTorch models into HLS code for FPGA synthesis, but it is not itself a search system.
% SNAC-Pack uses Optuna as its search backbone, using its shared SQLite storage for distributed parallel global search across compute nodes, and natively integrates hls4ml as the deployment target so that hardware metrics enter the optimization loop rather than being evaluated post-hoc.

% Recent work has also begun closing the loop between user intent and AutoML execution.
% AutoMMLab~\citep{yang2025autommlab} introduces a request-to-model workflow in which natural-language instructions drive an LLM agent through dataset selection, hyperparameter optimization, and deployment for computer vision tasks.
% SNAC-Pack adopts a related agentic interface through a Model Context Protocol layer that exposes its YAML-driven configuration, search, and synthesis stages as callable tools, lowering the barrier for non-experts who wish to target FPGA deployment on a new dataset.
Most AutoML stacks optimize software-side objectives (for example validation accuracy or throughput) and stop at a trained checkpoint~\citep{feurer2015efficient,jin2023autokeras,zimmer2021auto,wang2021flaml,tang2024autogluon}.
Optuna~\citep{optuna_2019} is a common programmatic backbone with support for distributed workers.
hls4ml translates quantized models to HLS for synthesis but is not a search system.
SNAC-Pack composes Optuna with shared SQLite-backed studies, rule4ml, and hls4ml so FPGA-relevant metrics enter the search loop instead of being evaluated only post hoc.
Related request-to-model agent lines such as AutoMMLab~\citep{yang2025autommlab} motivate our optional Model Context Protocol (MCP) front end, which exposes the same YAML-driven stages as tools for agentic clients; repository documentation lists the full tool surface.

\section{Method}
\label{sec:method}
SNAC-Pack builds on the two-stage global and local search workflow introduced in NAC~\citep{Weitz_2025}: a global stage explores architectural hyperparameters under multi-objective feedback, and a local stage compresses a small set of selected candidates with quantization-aware training (QAT) and pruning before synthesis with hls4ml.
SNAC-Pack extends NAC by implementing hardware surrogate objectives into the global loop (rule4ml~\citep{Rahimifar_2025}; wa-hls4ml predictors~\citep{hawks2025wahls4mlbenchmarksurrogatemodels} are discussed further in \S\ref{sec:related}) so that trial scores reflect FPGA-oriented costs without invoking full synthesis for every candidate.
The framework is designed as an end-to-end workflow for hardware-aware neural architecture search and FPGA deployment, where users configure datasets, search spaces, optimization objectives, hardware estimation settings, and deployment targets through a YAML interface.
Dataset loaders and splits, Optuna study settings, rule4ml/hls4ml-related options, and local-search schedules are part of the same configuration.
The workflow is depicted in Fig.~\ref{fig:pipeline_fig}; concrete hyperparameters for our experiments appear in the case studies so this section stays aligned with the implementation.

\subsection{Global Search}
\label{sec:method-global}
Global search runs multi-objective neural architecture search in Optuna~\citep{optuna_2019} with NSGA-II~\citep{NSGA}.
For each trial, the sampler draws architectural and training hyperparameters (depth, widths, activations, normalization, regularization, and other user-defined knobs) from the search space and builds a Keras model in TensorFlow.
The model is trained for a set number of epochs on the task loss, and a validation metric such as accuracy or readout fidelity is recorded.
When BOPs is an objective, it is computed analytically from the architecture and the configured numeric precision.
If hardware-aware mode is on, the same trained graph is passed to rule4ml's \texttt{MultiModelWrapper} together with an hls4ml-style configuration (board, strategy, precision, reuse factor, and related metadata) so LUT, FF, BRAM, DSP utilization and latency in clock cycles can be estimated quickly instead of invoking Vivado on every candidate.
Each trial returns one scalar per objective in a fixed order, and Optuna persists objectives and study metadata in a relational backend.

NSGA-II evolves a population under Pareto dominance, so the study keeps a set of mutually nondominated trials rather than collapsing everything to a single score.
In our setup, the task metric (accuracy for jet classification; readout fidelity for qubit readout, computed as validation assignment accuracy) is maximized while BOPs, surrogate resource metrics, and surrogate latency are minimized, which matches the deployment goal of keeping task quality high while pushing resource use and cycles down.
Which objectives are turned on is a YAML choice; in the experiments below we use subsets of the task metric, BOPs, mean predicted utilization across LUT/FF/BRAM/DSP, and rule4ml's latency field.

Rule4ml can score Keras or PyTorch models in general~\citep{Rahimifar_2025}; SNAC-Pack passes the trained Keras graph from each trial together with HLS settings to obtain predicted BRAM, DSP, FF, LUT counts or percentages and latency in clock cycles.
For search we aggregate the four utilization percentages into one average-resource objective (their arithmetic mean), while still logging the per-resource columns when we want to inspect constraints or post hoc behavior.
These numbers are inexpensive proxies: they help explore toward promising regions, while the hls4ml synthesis results after local search supply the ground truth we report in the case studies.

Trials may be run back-to-back on one machine or sent out to many workers or Slurm tasks that all point at the same Optuna storage URL.
The bundled workflow uses a shared SQLite study so workers can append completed trials safely; each worker executes the same code path as a single-node run, and ``parallel'' here means independent trials in flight at once, not model-parallel training inside one trial.
Wall time is dominated by short fine-tuning and surrogate calls, not by place-and-route.

\subsection{Local Search and Synthesis}
\label{sec:method-local}

Global search writes a Pareto archive to disk (CSV summaries and YAML checkpoints for selected trials).
SNAC-Pack picks one or more points on that front.
For example, best task performance subject to surrogate budgets is selected and exported as a YAML for the compression stage.

Local search uses SNAC-Pack's combined QAT\,+\,pruning driver (\texttt{tf\_local\_search\_combined}).
For each fixed-point precision pair $(\text{total bits}, \text{integer bits})$ in the local-search configuration, we convert the floating-point seed to a QKeras QAT model, run a QAT warmup for \texttt{qat\_epochs}, then run iterative magnitude pruning with TensorFlow Model Optimization's \texttt{prune\_low\_magnitude}.
Pruning repeats for a configured number of iterations; each pass removes another fraction of the weights that are still nonzero, with the aggressiveness set by the YAML \texttt{pruning\_rate} (concrete values appear in the case-study tables).
Each iteration trains the pruned QAT model for \texttt{epochs\_per\_iteration}, evaluates validation performance, strips pruning wrappers for measurement, logs sparsity and effective BOPs, and applies lottery-ticket-style weight rewinding: pruned weights are zeroed while surviving weights are reset to their post-warmup values before the next iteration.
The best weights seen across iterations for that precision are checkpointed, and the outer loop repeats over every precision pair so we obtain a small family of compressed models rather than a single point estimate.
Optional stratified $k$-fold splits reuse the same schedule per fold and average reported metrics; the fold count is part of the YAML experiment definition.

After local search, chosen checkpoints are exported through hls4ml using the board, I/O style, reuse factor, strategy, and fixed-point formats declared in the experiment configuration.
Vivado reports post-route resource utilization, latency, and initiation interval; those are the numbers we tabulate as ground truth in the case studies and use to check surrogate quality from global search.

\subsection{MCP Interface}
% To lower the barrier to entry for non-experts and enable agentic workflows, SNAC-Pack exposes its full pipeline through a Model Context Protocol (MCP) layer that wraps the search, planning, and synthesis stages as a fixed set of tools with typed schemas.
% The MCP server registers ten tools spanning dataset discovery (listing built-in datasets, profiling their modality, shape, and loader availability), search planning (deterministically mapping a dataset specification and user constraints to a complete YAML configuration, with explicit warnings when requests fall outside the supported model families), and end-to-end execution (running global search, local search, and synthesis from either an existing configuration or a freshly planned one).
% A natural-language entry point parses plain-English requests for search budget, hardware-aware objectives, latency or resource targets, and model-family preferences into structured planner constraints, allowing a user to run the full pipeline with a sentence rather than a configuration file.
% The same tool implementations are reachable through three integration paths: a native MCP server for MCP-capable clients, the OpenAI Responses API with remote MCP for managed tool orchestration, and a local tool-calling loop compatible with any OpenAI-style chat completions endpoint.
% All file paths passed to the tools are restricted to the repository root, so an agent cannot read or write outside the project.
% The underlying YAML-driven workflow is unchanged, and the MCP layer simply allows an LLM agent to inspect the repository, profile a new dataset, plan a hardware-aware search, and run it without writing code or configs by hand.
SNAC-Pack optionally exposes the same YAML-driven pipeline through a Model Context Protocol (MCP) layer so agents can plan searches, launch global and local search, and trigger synthesis without hand-editing configuration files.
Tool paths are restricted to the repository root to limit filesystem access; the README in the code repository summarizes the tool catalog, integration modes (native MCP, managed remote MCP, and local OpenAI-style tool calling), and natural-language entry points.

% \begin{table}[ht]
% \centering
% \caption{SNAC-Pack workflow}
% \label{tab:workflow}

% \begin{tabular}{ll}
% \toprule
% Step & Description \\
% \midrule

% 1 & Define YAML configuration for dataset, search space, objectives, and deployment settings \\

% 2 & Run global multi-objective architecture search \\

% 3 & Generate Pareto-optimal candidate architectures \\

% 4 & Run local search to refine selected models using QAT and pruning \\

% 5 & Synthesize models with hls4ml for FPGA deployment \\

% \bottomrule
% \end{tabular}

% \end{table}

\section{Case Studies}
\label{sec:case-studies}

\subsection{Jet Classification}

To show the effectiveness of SNAC-Pack, we apply it to jet classification, a common and challenging task in high energy physics at the Large Hadron Collider (LHC).
The goal is to accurately classify collision-created jets into one of five categories (light quark, gluon, W boson, Z boson, top quark) based on their kinematic properties.
This is showcased with the hls4ml LHC dataset \cite{pierini_2020_3602260}.
The 8 constituents with the greatest transverse momentum are used per jet.

For this task, we configure SNAC-Pack to search for an optimal multi-layer perceptron (MLP) architecture, with an 8 constituent MLP as a comparative baseline \cite{Odagiu_2024} with the data processed and normalized as done there.
This baseline is chosen, as it is one of the state-of-the-art architectures for this task.
% The global search uses the Non-dominated Sorting Genetic Algorithm II (NSGA-II) \cite{NSGA}, exploring a search space over depth, per-layer hidden widths, activations, batch normalization \cite{ioffe2015batch}, learning rate, $L_1$ regularization, and dropout, as listed in Table~\ref{tab:mlp-parameter-table} (Appendix~\ref{app:jet-search-config}).
% The objectives used are estimated average hardware utilization, estimated clock cycles, and accuracy. All training is performed with a batch size of 128. Global search is performed for 500 trials, 5 epochs per trial, and an evolutionary population size of 20.
% The search was ran on a single A100 GPU in serial: global search took approximately $6$\, hours, and the local search across all lsited bit precisions took approximately $8$\, hours.
Global search follows \S\ref{sec:method-global} (NSGA-II~\cite{NSGA}) over depth, per-layer widths, activations, batch normalization~\cite{ioffe2015batch}, learning rate, $L_1$ regularization, and dropout (Table~\ref{tab:mlp-parameter-table}, Appendix~\ref{app:jet-search-config}), maximizing accuracy while minimizing mean predicted utilization and surrogate latency.
We use batch size 128, 500 trials, five training epochs per trial, and population size 20, with three-fold cross-validation on one A100 in serial ($\sim$6\,h wall time for global search; $\sim$8\,h for local search across listed bit precisions).

The results are summarized in Fig.~\ref{fig:jet_pareto_combined}: three SNAC-Pack surrogate Pareto projections (average resources versus clock cycles, average resources versus accuracy, and clock cycles versus accuracy) and, for comparison, an NAC Pareto front in the BOPs-versus-accuracy plane.
For comparison, we also run NAC with the same number of trials and epochs, optimizing solely for BOPs and accuracy.
The resulting Pareto front is shown in Fig.~\ref{fig:nac_pareto_front}.
We take two checkpoints from the SNAC-Pack global Pareto archive to show that the workflow leaves multiple viable deployment choices rather than collapsing to a single scalarized solution.
The Optimal Accuracy SNAC-Pack row is the nondominated trial with highest validation accuracy subject to accuracy greater than 0.638, so it matches or exceeds the baseline~\cite{Odagiu_2024} on that axis.
The Optimal Resource SNAC-Pack row is the Pareto-optimal trial with the lowest BOPs, estimated average resource utilization, and estimated clock cycles among trials with validation accuracy at least 62\%.
Both architectures undergo the same local search and synthesis protocol (Table~\ref{tab:mlp-local-search-table}, Appendix~\ref{app:jet-search-config}) and appear alongside NAC and the baseline in Tables~\ref{tab:mlp-global-comparison} and~\ref{tab:model-comparison-synth}.

% For local search, each selected architecture is compressed with combined QAT and iterative magnitude pruning~\cite{quantization_survey,frankle2019stabilizing,frankle2018lottery}: a five-epoch QAT warm-up at 8-bit precision, followed by ten pruning iterations that each train for ten epochs and remove 20\% of the remaining weights on the quantized model.
% Local search and synthesis settings are summarized in Table~\ref{tab:mlp-local-search-table} (Appendix~\ref{app:jet-search-config}).
Local search follows \S\ref{sec:method-local} with combined QAT and magnitude pruning~\cite{quantization_survey,frankle2019stabilizing,frankle2018lottery}; concrete epochs, pruning rates, and board settings are in Table~\ref{tab:mlp-local-search-table} (Appendix~\ref{app:jet-search-config}).
With hls4ml, the resulting architectures are then synthesized on the Xilinx Virtex UltraScale+ VU13P FPGA, with io\_parallel io\_type, latency strategy, a reuse factor of 1.
Resource utilization and latency are shown in Table~\ref{tab:model-comparison-synth}.

\begin{table}[t]
    \caption{Comparison of model accuracy, BOPs, and estimated hardware metrics from global search.
    The Baseline was optimized for accuracy, NAC for accuracy and BOPs, and SNAC-Pack for accuracy, estimated average resources, and clock cycles.
    Two SNAC-Pack rows reflect different selection rules on the same Pareto archive.
    The best values are reported in bold.}
  \label{tab:mlp-global-comparison}
  \centering
  \footnotesize
  \begin{tabular}{lcccc}
    \toprule
    Model & Accuracy [\%] & BOPs & Est. average resources [\%]& Est. latency (cc) \\
    \midrule
    Baseline \cite{Odagiu_2024}          & 63.77 & 25{,}916 & 7.10 & 183.74 \\
    Optimal NAC \cite{Weitz_2025}      & 63.81 & 7{,}904  & 3.60   & 62.69 \\
    Optimal Accuracy SNAC-Pack & \textbf{63.84} & 8{,}352  & 3.12 & 72.24 \\
    Optimal Resource SNAC-Pack & 62.30 & \textbf{4{,}768} & \textbf{1.96} & \textbf{56.59} \\
    \bottomrule
  \end{tabular}
\end{table}

% Table~\ref{tab:mlp-global-comparison} shows that Optimal Accuracy SNAC-Pack reaches the highest validation accuracy in this comparison, while Optimal Resource SNAC-Pack trades about 1.5 percentage points of accuracy for substantially lower BOPs, estimated average resources, and estimated latency under its 62\% accuracy floor.
% After local search and synthesis (Table~\ref{tab:model-comparison-synth}), Optimal Resource SNAC-Pack improves post-route latency, initiation interval, LUT, FF, and BRAM relative to the baseline and the other learned models, demonstrating that the hardware-aware search objective can surface a genuinely cheaper deployment point when a modest accuracy relaxation is acceptable.
% Optimal Accuracy SNAC-Pack keeps the strongest accuracy among the searched models but ends up with higher post-synthesis latency than the baseline, which underscores that surrogate rankings during global search do not always align with Vivado timing and points to where the latency predictor is miscalibrated for this checkpoint.
Table~\ref{tab:mlp-global-comparison} summarizes surrogate scores at export; Table~\ref{tab:model-comparison-synth} gives post-route hardware.
Optimal Resource SNAC-Pack accepts a modest accuracy relaxation but delivers the best latency, II, LUT, FF, and BRAM among the searched rows, whereas Optimal Accuracy SNAC-Pack leads on validation accuracy yet is slower than the baseline after synthesis, illustrating surrogate--Vivado misalignment for latency on that checkpoint.

\begin{table}[t]
  \caption{Hardware resource utilization and latency for the selected models after local search and synthesis.
  The baseline is pruned by 50\% and quantized to 8 bits.
  NAC and both SNAC-Pack checkpoints are synthesized after the same style of local search. CC is the number of clock cycles.
  The SNAC-Pack architectures follow the pruning and QAT schedule in Table~\ref{tab:mlp-local-search-table}.
  The best values are reported in bold.}
  \label{tab:model-comparison-synth}
  \centering
  \scriptsize
  \resizebox{\textwidth}{!}{%
    \begin{tabular}{lcccccc}
      \toprule
      Model & Lat. [ns] (cc) & II [ns] (cc) & DSP & LUT & FF & BRAM \\
      \midrule
      Baseline \cite{Odagiu_2024} & 105 (21) & 5 (1) & 262 (2.1\%) & 155080 (9.0\%) & 25714 (0.7\%) & 4 (0.1\%) \\
      Optimal NAC \cite{Weitz_2025} & 125 (25) & 60 (12) & \textbf{0} & 54075 (3.13\%) & 12016 (0.35\%) & 8 (0.3\%) \\
      Optimal Accuracy SNAC-Pack & 140 (24) & 70 (12) & \textbf{0} & 57728 (3.34\%) & 12605 (0.36\%) & \textbf{0} \\
      Optimal Resource SNAC-Pack& \textbf{50 (10)} & \textbf{5 (1)} & \textbf{0} & \textbf{34472 (1.99\%)} & \textbf{5061 (0.15\%)} & \textbf{0} \\
      \bottomrule
    \end{tabular}%
  }
\end{table}

\begin{figure*}[t]
    \centering

    \begin{subfigure}{0.48\textwidth}
        \centering
        \includegraphics[width=\linewidth]{figures/old_results/aaa_avg_resource_vs_clock_cycles.pdf}
        \caption{Average resources vs. clock cycles.}
        \label{fig:snac-pack_res_cc}
    \end{subfigure}
    \hfill
    \begin{subfigure}{0.48\textwidth}
        \centering
        \includegraphics[width=\linewidth]{figures/old_results/aaa_val_acc_vs_avg_resource.pdf}
        \caption{Average resources vs. accuracy.}
        \label{fig:snac-pack_res_acc}
    \end{subfigure}

    \vspace{0.5em}

    \begin{subfigure}{0.48\textwidth}
        \centering
        \includegraphics[width=\linewidth]{figures/old_results/aaa_val_acc_vs_clock_cycles.pdf}
        \caption{Clock cycles vs. accuracy.}
        \label{fig:snac-pack_cc_acc}
    \end{subfigure}
    \hfill
    \begin{subfigure}{0.48\textwidth}
        \centering
        \includegraphics[width=\linewidth]{figures/old_results/aaa_val_acc_vs_bops.pdf}
        \caption{BOPs vs. accuracy (NAC).}
        \label{fig:nac_pareto_front}
    \end{subfigure}

    \caption{
    Pareto fronts obtained during jet classification global search. 
    SNAC-Pack optimizes estimated hardware-aware objectives and accuracy, while NAC optimizes BOPs and accuracy.
    Each point represents a sampled architecture.
    }
    \label{fig:jet_pareto_combined}
\end{figure*}

\subsection{Qubit Readout}

We also evaluate SNAC-Pack on qubit readout classification using superconducting qubit I/Q signal data~\cite{diguglielmo2025endtoendworkflowmachinelearningbased}.
The task is to classify the qubit state from time-series microwave readout signals under strict latency and hardware resource constraints.
The dataset consists of in-phase (I) and quadrature (Q) signals collected from dispersive transmon qubit measurements, where the measured microwave response depends on the underlying qubit state.
Readout fidelity is the validation assignment accuracy (\%): the fraction of readout windows whose predicted binary qubit state matches the label.

% For this task, we configure SNAC-Pack to search for compact MLP-based architectures optimized for FPGA deployment using hls4ml.
% The global search uses Optuna with the Non-dominated Sorting Genetic Algorithm II (NSGA-II)~\cite{NSGA}, exploring hidden widths, activations, normalization, and readout-window settings summarized in Table~\ref{tab:qubit-global-parameter-table} (Appendix~\ref{app:qubit-search-config}).
% The study optimizes four objectives: maximize validation fidelity, minimize normalized estimated FPGA resource utilization from rule4ml, estimated latency in clock cycles, and BOPs.
% We distribute the workload over ten Slurm compute nodes, each allocated one NVIDIA A100 GPU and running fifty trials, for five hundred trials in total; each trial trains for five epochs with stratified three-fold cross-validation to obtain the objective values, using a batch size of 128.
% Global search took approximately $3$\,hours per node, with all ten nodes running in parallel; local search took approximately $1.5$\,hours per node, with the four fixed-point precisions each run on its own node in parallel.
% When global search finishes, we take the completed trial with the highest validation fidelity and export that architecture for the subsequent compression stage, using the default filename from the tutorial scripts (\texttt{best\_model\_for\_local\_search.yaml}).
% Local search applies combined quantization-aware training and iterative magnitude pruning under the schedule and fixed-point sweeps recorded in Table~\ref{tab:qubit-local-search-table} (Appendix~\ref{app:qubit-search-config}), including twenty warm-up epochs, ten pruning rounds with twenty training epochs per round, a pruning-rate parameter of $20\%$ per iteration, and four decreasing precision pairs unless a run-specific precision override is set via the environment, then compilation and synthesis with hls4ml and Vivado as reported in Table~\ref{tab:qubit-model-comparison-synth}.
For this task, we configure SNAC-Pack to search over compact MLP-based architectures for hls4ml deployment.
Global search follows \S\ref{sec:method-global} with NSGA-II~\cite{NSGA} in Optuna over Table~\ref{tab:qubit-global-parameter-table} (Appendix~\ref{app:qubit-search-config}), maximizing readout fidelity while minimizing rule4ml resource utilization, clock cycles, and BOPs.
We distribute 500 trials (fifty per node) across ten Slurm nodes with one NVIDIA A100 GPU each, five epochs per trial, stratified three-fold cross-validation, and batch size 128 ($\sim$3\,h per node wall time for global search in parallel; $\sim$1.5\,h per node for local search with four precisions across nodes).
We export the highest-fidelity trial as \texttt{best\_model\_for\_local\_search.yaml} for local search under Table~\ref{tab:qubit-local-search-table} (Appendix~\ref{app:qubit-search-config}), then compiled and synthesized with  hls4ml and Vivado as reported in Table~\ref{tab:qubit-model-comparison-synth}.

The baseline model uses a 400 clock cycle readout window starting at clock cycle 100, where the I and Q signal windows are concatenated into an 800-dimensional input.
SNAC-Pack explores smaller readout windows and reduced model capacity to reduce latency and hardware utilization while maintaining readout fidelity.
The search also allows optional removal of hidden activations to evaluate whether nonlinearities improve performance enough to justify additional deployment cost.

Figure~\ref{fig:qubit_pareto_combined} summarizes surrogate trade-offs from the completed study: each panel projects the sampled architectures in the plane of validation fidelity against one hardware-aware axis (average resources, clock cycles, or BOPs), so the cloud of points traces how much fidelity must be relaxed as estimated utilization, latency, or operator count is pushed down.
Unlike the jet experiment, we report a single SNAC-Pack checkpoint that matches the export rule above. The trial with highest validation fidelity is taken from the archive and carried through the same local search and synthesis protocol, which appears alongside the baseline in Tables~\ref{tab:qubit-global-comparison} and~\ref{tab:qubit-model-comparison-synth}.

\begin{table}[t]
  \caption{Comparison of readout fidelity, BOPs, and estimated hardware metrics from global search.
  Readout fidelity denotes validation assignment accuracy (\%) on the scored split.
  The baseline was tuned for fidelity alone; SNAC-Pack was searched for fidelity together with estimated average resources and clock cycles.
  The best values are reported in bold.}
  \label{tab:qubit-global-comparison}
  \centering
  \footnotesize
  \begin{tabular}{lcccc}
    \toprule
    Model & Fidelity [\%] & BOPs & Est. average resources [\%] & Est. latency (cc) \\
    \midrule
    Baseline \cite{diguglielmo2025endtoendworkflowmachinelearningbased} & \textbf{96.06} & 3{,}487{,}250 & 0.79 & 948.67 \\
    Optimal SNAC-Pack & 95.21 & \textbf{1{,}026{,}048} & \textbf{0.61} & \textbf{810.36} \\
    \bottomrule
  \end{tabular}
\end{table}

\begin{table}[t]
  \caption{Hardware resource utilization and latency estimates for qubit readout models synthesized with hls4ml.
  The baseline model uses a 400 clock cycle readout window with an 800-dimensional input.
  The SNAC-Pack model is synthesized after local search with iterative pruning and quantization-aware training.
  CC denotes the number of clock cycles.
  The best values are reported in bold.}
  \label{tab:qubit-model-comparison-synth}
  \centering
  \scriptsize
  \resizebox{\textwidth}{!}{%
    \begin{tabular}{lcccccc}
      \toprule
      Model & Lat. [ns] (cc) & II [ns] (cc) & DSP & LUT & FF & BRAM \\
      \midrule
      Baseline & 19.35 (6) & 16.1 (5) & \textbf{0} & 15045 (5.49\%) & 3656 (0.67\%) & \textbf{0} \\
      Optimal SNAC-Pack & \textbf{12.9 (4)} & \textbf{12.9 (4)} & \textbf{0} & \textbf{6996 (2.55\%)} & \textbf{2330 (0.43\%)} & \textbf{0} \\
      \bottomrule
    \end{tabular}%
  }
\end{table}

Table~\ref{tab:qubit-global-comparison} makes the surrogate trade-off explicit for that exported design: relative to the baseline tuned for fidelity alone, readout fidelity moves from $96.06\%$ to $95.21\%$ while BOPs fall from $3{,}487{,}250$ to $1{,}026{,}048$, estimated average resource utilization drops from $0.79\%$ to $0.61\%$, and the rule4ml clock-cycle latency estimate decreases from $948.67$ to $810.36$.
The task-metric loss is modest compared with the reduction in operator count and the surrogate hardware scores, consistent with the broad achievable region illustrated in Fig.~\ref{fig:qubit_pareto_combined}.

After local search and Vivado synthesis, Table~\ref{tab:qubit-model-comparison-synth} shows that the same checkpoint also occupies a cheaper post-route point than the baseline: latency and initiation interval shrink to $12.9\,\mathrm{ns}$ (four clock cycles) from $19.35\,\mathrm{ns}$ (six cycles), while LUT and FF counts are reduced to roughly half of the baseline footprint (DSP and BRAM remain at zero in both rows).
Because the search varies readout windowing and layer sizes, these results additionally indicate that competitive fidelity can be maintained away from the fixed $400$-cycle baseline window while still improving on-chip timing and area.

\begin{figure*}[t]
    \centering

    \begin{subfigure}{0.48\textwidth}
        \centering
        \includegraphics[width=\linewidth]{figures/qubit_optuna_job_52893384_01_avg_resource_vs_clock_cycles.pdf}
        \caption{Average resources vs. clock cycles.}
        \label{fig:snac-pack_qubit_res_cc}
    \end{subfigure}
    \hfill
    \begin{subfigure}{0.48\textwidth}
        \centering
        \includegraphics[width=\linewidth]{figures/qubit_optuna_job_52893384_02_avg_resource_vs_performance_metric.pdf}
        \caption{Average resources vs. accuracy.}
        \label{fig:snac-pack_qubit_res_acc}
    \end{subfigure}

    \vspace{0.5em}

    \begin{subfigure}{0.48\textwidth}
        \centering
        \includegraphics[width=\linewidth]{figures/qubit_optuna_job_52893384_03_clock_cycles_vs_performance_metric.pdf}
        \caption{Clock cycles vs. accuracy.}
        \label{fig:snac-pack_qubit_cc_acc}
    \end{subfigure}
    \hfill
    \begin{subfigure}{0.48\textwidth}
        \centering
        \includegraphics[width=\linewidth]{figures/qubit_optuna_job_52893384_04_bops_vs_performance_metric.pdf}
        \caption{BOPs vs. accuracy.}
        \label{fig:snac-pack_qubit_bops_acc}
    \end{subfigure}

    \caption{
    Pareto fronts obtained during qubit readout global search. 
    SNAC-Pack optimizes estimated hardware-aware objectives, readout fidelity, and BOPs.
    Each point represents a sampled architecture.
    }
    \label{fig:qubit_pareto_combined}
\end{figure*}

\section{Conclusion}
This work introduced SNAC-Pack, an open-source AutoML pipeline for hardware-aware neural architecture search, model compression, and FPGA deployment.
Built on Optuna and parallel trial workers, SNAC-Pack extends NAC by scoring global search trials with rule4ml surrogate estimates of utilization and latency.

On LHC jet classification, NSGA-II produces a Pareto set of models instead of a single best design. The resource-focused SNAC-Pack model does better than the baseline and NAC on post-route LUT, flip-flop, BRAM, latency, and initiation interval after the same compression setup (Table~\ref{tab:model-comparison-synth}). The accuracy-focused model stays strong on validation accuracy but is slower than the baseline after synthesis, which shows that surrogate rank during search can disagree with Vivado timing and that predictors can be improved.

On superconducting qubit readout, we export the trial with highest validation fidelity and run the same QAT and pruning pipeline. Readout fidelity drops slightly compared to the baseline while BOPs and surrogate hardware scores improve, and post-route latency and area on the ZCU102 beat the tabulated baseline (Tables~\ref{tab:qubit-global-comparison} and~\ref{tab:qubit-model-comparison-synth}).

Users configure datasets, objectives, and hardware settings through YAML and an optional MCP layer allows for the use of the same sandboxed tools.
Future work includes improving and calibrating surrogate estimation, expanding search spaces as predictors support them, and reporting search variance across seeds and splits.

\clearpage
\newpage
\bibliography{references}

\newpage
\section{[Optional] Submission Checklist}
% All submissions are welcome to include the AutoML submission checklist provided
% below. \textbf{While not required, including the checklist may be beneficial
%   during the reproducibility review process and for maximizing long-term
%   impact.} If you choose to include the checklist, it will not count towards the
% page limit. The submission checklist draws upon related submission checklists:
% %
% \href{https://neurips.cc/Conferences/2022/PaperInformation/PaperChecklist}
%      {the NeurIPS '22 checklist}
% %
% and the
% \href{https://www.automl.org/wp-content/uploads/NAS/NAS_checklist.pdf}
%      {\textsc{nas} checklist}.
% %

% For each question, change the default \verb|\answerTODO{}| (typeset \answerTODO)
% to
% \verb|\answerYes{[justification]}| (typeset \answerYes),
% \verb|\answerNo{[justification]}| (typeset \answerNo), or
% \verb|\answerNA{[justification]}| (typeset \answerNA).
% We recommend including a brief justification to your answer, either by
% referencing the appropriate section of your paper or providing a brief inline
% description.  For example:
% \begin{itemize}
% \item Did you include the license of the code and datasets? \answerYes{See
%   Section~\ref{sec:code}.}
% \item Did you include all the code for running experiments? \answerNo{We include
%   the code we wrote for conducting the experiments, but complete replication
%   depends on proprietary libraries for executing on a private compute
%   cluster. The code therefore is not runnable without modification. To
%   compensate, we provide a runnable but non-parallelized version of the code
%   that could replicate the results at the expense of a greater wall-clock time.}
% \item Did you include the license of the datasets? \answerNA{Our experiments
%   were conducted on publicly available datasets and we have not introduced new
%   datasets.}
% \end{itemize}
% Please note that if you answer a question with \verb|\answerNo{}|, we recommend
% providing an explanation and/or compensation for the omission. For example, if
% you cannot provide complete evaluation code for some reason, you might instead
% provide code for a minimal reproduction of the main insights of your paper.

% Please do not modify the questions and only use the provided macros for your
% answers. Note that the submission checklist does not count towards the page
% limit. If you choose to modify \texttt{instructions.tex}, please delete these
% instructions and only keep the Submission Checklist section heading above along
% with the questions/answers below.

\begin{enumerate}
\item For all authors\dots
  %
  \begin{enumerate}
  \item Do the main claims made in the abstract and introduction accurately
    reflect the paper's contributions and scope?
    %
    \answerYes{The abstract and Section~\ref{sec:intro} describe SNAC-Pack as an open-source hardware-aware NAS framework with surrogate resource estimation, distributed search, and end-to-end FPGA deployment, all of which are demonstrated in the case studies (Section~\ref{sec:case-studies}).}
    %
  \item Did you describe the limitations of your work?
    %
    \answerYes{We discuss limitations in the Conclusion, including the accuracy of the rule4ml surrogate (Table~\ref{tab:model-comparison-synth} shows the SNAC-Pack model has slightly higher latency than predicted), the restriction to MLP search spaces in the present case studies, and the lack of robustness/fault-tolerance objectives.}
    %
  \item Did you discuss any potential negative societal impacts of your work?
    %
    \answerYes{SNAC-Pack is a general-purpose AutoML framework for FPGA deployment evaluated on scientific ML tasks (LHC jet classification and qubit readout). The main consideration specific to this work is the agentic MCP interface, which exposes the pipeline to LLM-driven tool calls; to mitigate misuse and accidental filesystem access, all path arguments are confined to the repository root, and tool implementations are thin wrappers over the same YAML-driven workflow used in the non-agentic mode. Beyond this, we do not foresee direct negative societal impacts beyond the standard concerns about automating ML pipelines.}
    %
  \item Did you read the ethics review guidelines and ensure that your paper
    conforms to them? (see \url{https://2022.automl.cc/ethics-accessibility/})
    %
    \answerYes{}
    %
  \end{enumerate}
  %
\item If you ran experiments\dots
  %
  \begin{enumerate}
  \item Did you use the same evaluation protocol for all methods being compared (e.g.,
    same benchmarks, data (sub)sets, available resources, etc.)?
    %
    \answerYes{NAC and SNAC-Pack are compared on the same hls4ml LHC jet dataset~\citep{pierini_2020_3602260} with the same number of trials (500), epochs per trial (5), and evolutionary population size (20), and synthesized on the same VU13P FPGA with identical hls4ml settings (Table~\ref{tab:mlp-local-search-table}).}
    %
  \item Did you specify all the necessary details of your evaluation (e.g., data splits,
    pre-processing, search spaces, hyperparameter tuning details and results, etc.)?
    %
    \answerYes{Search spaces are given in Tables~\ref{tab:mlp-parameter-table} and~\ref{tab:qubit-global-parameter-table}; local search and synthesis settings are in Tables~\ref{tab:mlp-local-search-table} and~\ref{tab:qubit-local-search-table}; data preprocessing follows~\citep{Odagiu_2024}.}
    %
  \item Did you repeat your experiments (e.g., across multiple random seeds or
    splits) to account for the impact of randomness in your methods or data?
    %
    \answerYes{We ran the searches with k-fold cross-validation with random seeds to ensure statistical robustness.}
    %
  \item Did you report the uncertainty of your results (e.g., the standard error
    across random seeds or splits)?
    %
    \answerNo{Accuracy, BOPs, and post-synthesis resource numbers are reported as point estimates from a single run. See the answer to (c).}
    %
  \item Did you report the statistical significance of your results?
    %
    \answerNo{Not reported, for the same reason as (d).}
    %
  \item Did you use enough repetitions, datasets, and/or benchmarks to support
    your claims?
    %
    \answerYes{We evaluate SNAC-Pack on two distinct scientific ML tasks (LHC jet classification and superconducting qubit readout) with different input modalities, search spaces, and target FPGAs (VU13P and ZCU102), demonstrating that the framework generalizes across deployment settings.}
    %
  \item Did you compare performance over time and describe how you selected the
    maximum runtime?
    %
    \answerNo{We report total trial counts (500 trials, 5 epochs each for global search). The runtime budget was selected to match the NAC comparison setup~\citep{Weitz_2025}.}
    %
  \item Did you include the total amount of compute and the type of resources
    used (e.g., type of \textsc{gpu}s, internal cluster, or cloud provider)?
    %
    \answerYes{See Section~\ref{sec:case-studies}}
    %
  \item Did you run ablation studies to assess the impact of different
    components of your approach?
    %
    \answerYes{The comparison against NAC~\citep{Weitz_2025} in Tables~\ref{tab:mlp-global-comparison} and~\ref{tab:model-comparison-synth} serves as an ablation isolating the contribution of the rule4ml surrogate objectives versus the BOPs proxy under an otherwise identical pipeline.}
    %
  \end{enumerate}
  %
\item With respect to the code used to obtain your results\dots
  %
  \begin{enumerate}
\item Did you include the code, data, and instructions needed to reproduce the
    main experimental results, including all dependencies (e.g.,
    \texttt{requirements.txt} with explicit versions), random seeds, an instructive
    \texttt{README} with installation instructions, and execution commands
    (either in the supplemental material or as a \textsc{url})?
    %
    \answerYes{The framework is open-source at \url{https://github.com/fastmachinelearning/nac-opt}, including the YAML configurations for both case studies, environment specification, and a README with installation and execution instructions.}
    %
  \item Did you include a minimal example to replicate results on a small subset
    of the experiments or on toy data?
    %
    \answerYes{The repository ships tutorials for the jet classification and qubit readout case studies as runnable examples.}
    %
  \item Did you ensure sufficient code quality and documentation so that someone else
    can execute and understand your code?
    %
    \answerYes{The repository includes per-module documentation. The MLP and block search paths share a unified interface, and the configuration is YAML-driven so external users can target new datasets without modifying source.}
    %
  \item Did you include the raw results of running your experiments with the given
    code, data, and instructions?
    %
    \answerYes{Pareto front data underlying Figs.~\ref{fig:snac-pack_res_cc}–\ref{fig:snac-pack_qubit_bops_acc} are saved as Optuna SQLite study files in the repository.}
    %
  \item Did you include the code, additional data, and instructions needed to generate
    the figures and tables in your paper based on the raw results?
    %
    \answerYes{Plotting scripts for the Pareto fronts are included in the repository.}
    %
  \end{enumerate}
  %
\item If you used existing assets (e.g., code, data, models)\dots
  %
  \begin{enumerate}
  \item Did you cite the creators of used assets?
    %
    \answerYes{We cite the hls4ml LHC jet dataset~\citep{pierini_2020_3602260}, the qubit readout dataset~\citep{diguglielmo2025endtoendworkflowmachinelearningbased}, NAC~\citep{Weitz_2025}, rule4ml~\citep{Rahimifar_2025}, hls4ml~\citep{hls4ml,fastml_hls4ml}, and Optuna~\citep{optuna_2019}.}
    %
  \item Did you discuss whether and how consent was obtained from people whose
    data you're using/curating if the license requires it?
    %
    \answerNA{All datasets used are publicly released for ML research.}
    %
  \item Did you discuss whether the data you are using/curating contains
    personally identifiable information or offensive content?
    %
    \answerNA{The datasets contain particle physics simulation data and superconducting qubit measurements, with no personally identifiable or offensive content.}
    %
  \end{enumerate}
  %
\item If you created/released new assets (e.g., code, data, models)\dots
  %
  \begin{enumerate}
    %
    \item Did you mention the license of the new assets (e.g., as part of your
    code submission)?
    %
    \answerYes{Apache-2.0 license is mentioned in the repository. See below.}
    %
    \item Did you include the new assets either in the supplemental material or as
    a \textsc{url} (to, e.g., GitHub or Hugging Face)?
    %
    \answerYes{\url{https://github.com/fastmachinelearning/nac-opt}}
    %
  \end{enumerate}
  %
\item If you used crowdsourcing or conducted research with human subjects\dots
  %
  \begin{enumerate}
  \item Did you include the full text of instructions given to participants and
    screenshots, if applicable?
    %
    \answerNA{No crowdsourcing or human subjects research was conducted.}
    %
  \item Did you describe any potential participant risks, with links to
    institutional review board (\textsc{irb}) approvals, if applicable?
    %
    \answerNA{No crowdsourcing or human subjects research was conducted.}
    %
  \item Did you include the estimated hourly wage paid to participants and the
    total amount spent on participant compensation?
    %
    \answerNA{No crowdsourcing or human subjects research was conducted.}
    %
  \end{enumerate}
\item If you included theoretical results\dots
  %
  \begin{enumerate}
  \item Did you state the full set of assumptions of all theoretical results?
    %
    \answerNA{This paper does not present theoretical results.}
    %
  \item Did you include complete proofs of all theoretical results?
    %
    \answerNA{This paper does not present theoretical results.}
    %
  \end{enumerate}
  %
\end{enumerate}

\newpage
\appendix
\section{Supplemental Section}

\subsection{Jet classification search space and configuration}
\label{app:jet-search-config}

\begin{table}[H]
  \caption{Comprehensive parameter space for jet classification.}
  \label{tab:mlp-parameter-table}
  \centering
  \small
  \renewcommand{\arraystretch}{1.2}
  \begin{tabular}{ll}
    \toprule
    \textbf{Parameter} & \textbf{Space} \\
    \midrule
    Number of layers        & \{4, 5, 6, 7, 8\} \\
    Hidden units per layer  & \\
    \quad Layer 1 & \{64, 120, 128\} \\
    \quad Layer 2 & \{32, 60, 64\} \\
    \quad Layer 3 & \{16, 32\} \\
    \quad Layer 4 & \{32, 64\} \\
    \quad Layer 5 & \{32, 64\} \\
    \quad Layer 6 & \{32, 64\} \\
    \quad Layer 7 & \{16, 32\} \\
    \quad Layer 8 & \{32, 44, 64\} \\
    Activation function     & \{ReLU, Tanh, Sigmoid\} \\
    Batch normalization     & \{True, False\} \\
    Learning rate           & \{0.0010, 0.0015, 0.0020\} \\
    L1 regularization       & \{0.0, $10^{-6}$, $10^{-5}$, $10^{-4}$\} \\
    Dropout rate            & \{0.0, 0.05, 0.1\} \\
    \bottomrule
  \end{tabular}
\end{table}

\begin{table}[H]
  \caption{Local search and synthesis configuration for jet classification.}
  \label{tab:mlp-local-search-table}
  \centering
  \small
  \renewcommand{\arraystretch}{1.2}
  \begin{tabular}{ll}
    \toprule
    \textbf{Parameter} & \textbf{Value / Space} \\
    \midrule
    Warm-up epochs & 5 \\
    Pruning iterations & 10 \\
    Pruning epochs per iteration & 10 \\
    Pruning rate per iteration & 20\% \\
    Quantization precision & \{(32,16), (16,6), (8,3), (4,1)\} \\
    FPGA board & VU13P \\
    HLS io\_type & io\_parallel \\
    HLS strategy & Latency \\
    Reuse factor & 1 \\
    \bottomrule
  \end{tabular}
\end{table}

\subsection{Qubit readout search space and configuration}
\label{app:qubit-search-config}

\begin{table}[H]
  \caption{Global search parameter space for qubit readout classification.}
  \label{tab:qubit-global-parameter-table}
  \centering
  \small
  \renewcommand{\arraystretch}{1.2}
  \begin{tabular}{ll}
    \toprule
    \textbf{Parameter} & \textbf{Space} \\
    \midrule
    Hidden units & \{2, 4\} \\
    Hidden activation & \{ReLU, LeakyReLU, None\} \\
    Normalization & \{BatchNorm, None\} \\
    Block type & \{None\} \\
    Number of blocks & \{1\} \\
    MLP head layers & \{2\} \\
    Output dimension & \{1\} (binary readout; scalar labels) \\
    Output activation & \{None (logits; binary cross-entropy)\} \\
    Readout window size &
    \{25, 50, 75, ..., 400\} \\
    Window start location &
    \{0, 25, 50, ..., 725\} \\
    \bottomrule
  \end{tabular}
\end{table}

\begin{table}[H]
  \caption{Local search and synthesis configuration for qubit readout classification.}
  \label{tab:qubit-local-search-table}
  \centering
  \small
  \renewcommand{\arraystretch}{1.2}
  \begin{tabular}{ll}
    \toprule
    \textbf{Parameter} & \textbf{Value / Space} \\
    \midrule
    Warm-up epochs & 20 \\
    Pruning iterations & 10 \\
    Pruning epochs per iteration & 20 \\
    Pruning rate per iteration & 20\% \\
    Quantization precision & \{(32,16), (16,6), (8,3), (4,1)\} \\
    FPGA board & ZCU102 \\
    HLS strategy & Latency \\
    Reuse factor & 8 \\
    \bottomrule
  \end{tabular}
\end{table}

\end{document}
